\section{Related Work}
\label{sec:related}

\paragraph{Video-generation evaluation.}
General-purpose suites assess visual quality, temporal consistency, motion, and semantic alignment. VBench~\cite{huang2023vbench} decomposes quality into dimensions including background consistency, motion smoothness and dynamic degree, and VBench-2.0~\cite{Zheng2025VBench20AV} extends it toward intrinsic faithfulness; EvalCrafter~\cite{eval_crafter} learns a weighting of many metrics from user opinion; Video-Bench~\cite{han2025videobench} scores generations with multimodal language models validated against human preference. Others specialise: T2V-CompBench~\cite{sun2024t2v_compbench} targets compositional binding, VMBench~\cite{ling2025vmbench} perceptually aligned motion quality, and WorldScore~\cite{duan2025worldscore} controllability and dynamics for world generation. These set the standards a benchmark paper must meet---explicit taxonomy, validated metrics, stated scale---but their measurements are whole-frame or sequence-level, and such an aggregate cannot reveal whether motion originates where intended or in displacement of nominally static structure. VMBench is the closest comparator, but it asks whether motion is perceptually plausible where we ask where it occurs and whether it survives compensation for global displacement. Our validation is correspondingly mechanistic---perturb by a known amount, require the expected response---complementing the human-alignment evidence those suites provide.

\paragraph{Long-horizon autoregressive generation.}
Autoregressive and streaming generators extend synthesis beyond short bidirectional clips by producing frames or temporal chunks causally, including CausVid~\cite{yin2025causvid}, Self-Forcing~\cite{huang2025selfforcing}, Infinite-Forcing~\cite{infinite-forcing}, Rolling-Forcing~\cite{liu2025rolling}, Reward-Forcing~\cite{lu2025reward}, LongLive~\cite{yang2025longlive}, and Causal-Forcing~\cite{zhu2026causal} with its few-step successor Causal-Forcing++~\cite{zhao2026causalforcingpp}; bidirectional systems such as LTX-Video~\cite{hacohen2024ltxvideo} and Wan~\cite{wan2025} provide short-horizon reference points. These systems differ substantially in temporal block size, training objective, memory and cache policy, and horizon-extension mechanism~\cite{yesiltepe2025infinity}, yet their evaluation typically relies on generic video metrics or method-specific qualitative analysis. That diversity motivates auditing each system in its released configuration, and the evaluation gap motivates a system-agnostic question: does static support remain fixed while dynamic flow persists?

\paragraph{Image dynamics and nature-flow animation.}
Work on animating still images has long modelled natural motion explicitly: controllable fluid animation separates static image structure from specified fluid motion~\cite{mahapatra2022controllable}, and Generative Image Dynamics models long-term scene motion through learned trajectory representations~\cite{Li_2024_gen_image_dynamics}. These works treat static support and dynamic content as distinct modelling targets; SNF-Bench elevates that distinction to an evaluation principle for long-horizon generators.
